\section{Conclusão}

A consolidação computacional confirmou as expectativas teóricas estabelecidas
nas seções iniciais. A implementação simbólica serviu como guia para validar
a correta geração dos termos e possibilitou comparar diretamente os resultados
com as versões otimizadas. O uso de matrizes de projeto vetorizadas e a
decomposição em blocos reduziu o tempo de treinamento em duas ordens de
grandeza, abrindo espaço para experimentos com milhares de amostras e graus
médios. Por fim, os experimentos demonstraram que o ganho de fidelidade com o
aumento de $r$ apresenta retornos decrescentes, reforçando a necessidade de
critérios adaptativos ou regularização adicional em aplicações de maior escala.

Como continuidade direta, pretendemos explorar outros conjuntos de funções base
- como polinômios ortogonais, funções radiais, splines e wavelets - além de
investigar estratégias de seleção automática de grau polinomial e
regularização adaptativa para cenários com ruído não gaussiano. Outro
desdobramento relevante é generalizar a estrutura da regressão para acomodar
funções alvo com características heterogêneas, incluindo dados tensoriais e
categorias discretas. Essas extensões visam ampliar a aplicabilidade do método
em diferentes domínios, como aprendizado de dinâmica de sistemas, modelagem
física e análise de dados multidimensionais.
